\section{A proof-oriented AESP-to-LocSOR hybrid}
\label{sec:hybrid-algorithm}

\subsection{Design principle}

The proposed division of labor is
\begin{equation}
    \boxed{
    \text{AESP for early accelerated compression}
    \quad\longrightarrow\quad
    \text{LocSOR for local high-accuracy cleanup}.}
    \label{eq:hybrid-principle}
\end{equation}
AESP is used only while its shifted systems and small early active sets are
expected to provide favorable progress per edge access.  The tail discards all
outer momentum and runs a direct local coordinate method on the original
PageRank objective.  The proof-oriented version uses \(\omega=1\) in the tail;
a safeguarded over-relaxed variant is developed in \cref{sec:sor-correction}.

The clean theoretical handoff is based on the original objective gap.  Define
\begin{equation}
    \Delta_{\rm sw}
    :=\frac{\alpha^{3/2}\epsppr}{1+\alpha}.
    \label{eq:switch-gap}
\end{equation}
AESP supplies the computable upper bound \(U_t\) in
\cref{eq:aesp-objective-bound}.  The switch occurs at
\begin{equation}
    J
    :=\min\{t\geq0:U_t\leq\Delta_{\rm sw}\}.
    \label{eq:switch-stage}
\end{equation}
This rule is conservative but does not require knowing \(f(x^0)\).

\subsection{Algorithm}

\begin{algorithm}[t]
\caption{Proof-oriented hybrid AESP--LocSOR}
\label{alg:hybrid-aesp-locsor}
\begin{algorithmic}[1]
\REQUIRE Graph \(G\), seed distribution \(s=e_v\),
         \(\alpha\in(0,1/2)\), target \(\epsppr>0\)
\STATE Set \(x^{(0)}=y^{(0)}=0\), \(\kappa_{\mathrm A}=1-2\alpha\), and define \(U_t\) by
       \cref{eq:aesp-objective-bound}
\FOR{\(t=1,2,\ldots\)}
    \STATE Approximately minimize
    \(
       h_t(z)=f(z)+\frac{\kappa_{\mathrm A}}{2}\norm{z-y^{(t-1)}}_2^2
    \)
    with LOCAPPR or LOCGD to the AESP-prescribed accuracy \(\phi_t\), obtaining
    \(x^{(t)}\)
    \IF{\(\norm{D^{-1/2}\nabla f(x^{(t)})}_\infty<\alpha\epsppr\)}
        \RETURN \(D^{1/2}x^{(t)}\)
    \ENDIF
    \IF{\(U_t\leq\Delta_{\rm sw}\)}
        \STATE Set \(x_{\rm sw}=x^{(t)}\) and break
    \ENDIF
    \STATE Update
    \(
       y^{(t)}=x^{(t)}+\beta_{\mathrm A}(x^{(t)}-x^{(t-1)})
    \)
    with \(\beta_{\mathrm A}\) from \cref{eq:aesp-momentum}
\ENDFOR
\STATE Discard all Catalyst momentum and initialize
       \(x\leftarrow x_{\rm sw}\), \(\nabla f(x)\leftarrow Qx-b\)
\WHILE{there exists \(u\) with
       \(\abs{\nabla_u f(x)}\geq\alpha\epsppr\sqrt{d_u}\)}
    \STATE Choose any active \(u\) and set
    \(
       x\leftarrow x-\frac{2\nabla_u f(x)}{1+\alpha}e_u
    \)
    \STATE Update \(\nabla f(x)\leftarrow Qx-b\) locally on
          \(\{u\}\cup\mathcal N(u)\)
\ENDWHILE
\RETURN \(D^{1/2}x\)
\end{algorithmic}
\end{algorithm}

\subsection{State transfer and invariants}

Three details are mandatory.

\paragraph{Use the original gradient.}
The inner queue is controlled by
\begin{equation}
    \nabla h_t(x)
    =\nabla f(x)+\kappa_{\mathrm A}(x-y^{(t-1)}).
    \label{eq:inner-vs-original-gradient}
\end{equation}
The tail must instead initialize from \(\nabla f(x_{\rm sw})\).  Using the inner
proximal gradient would solve the wrong linear system.

\paragraph{Reset momentum.}
The vectors \(x^{(t)}-x^{(t-1)}\) in Catalyst and the momentum state used by
Chebyshev or heavy-ball iterations arise from different recurrences.  The
tail therefore starts with zero momentum.  Reusing Catalyst momentum in a
LocCH tail would not be justified by either source analysis.

\paragraph{Preserve the final certificate.}
The tail stops only when \cref{eq:ppr-certificate} holds.  An inner AESP
threshold \(\epsin^{(t)}\) certifies the shifted subproblem and
is not a substitute for the original PPR residual test.

\subsection{Practical switching rules}

The theoretical rule \cref{eq:switch-stage} may be conservative.  The
following observable rules are reasonable experimental variants, but they do
not replace the theorem unless separately analyzed.

\begin{enumerate}[leftmargin=2em]
\item \textbf{Residual-mass switch.}  Maintain the best checkpoint in
\(
M_t=\norm{D^{1/2}\nabla f(x^{(t)})}_1
\)
and switch when \(M_t\leq\theta\alpha^{3/2}\).  This gives the stronger
LocSOR bounds in \cref{sec:mass-handoff}, but the AESP burn-in time needed to
reach this threshold is currently open.

\item \textbf{Marginal-work switch.}  Let \(W_t\) be the work of the latest
AESP stage and let \(P_t\) be an observable progress statistic, such as the
reduction of the best original residual mass.  Switch when \(P_t/W_t\) remains
below a calibrated threshold for several stages.

\item \textbf{Probe switch.}  Checkpoint the AESP state, run a short LocSOR or
LocCH probe, and accept the handoff only if the probe provides sufficient
certified progress per work.  Otherwise roll back and continue AESP.
\end{enumerate}

The benchmark in \cref{sec:tbscience} should leave the switching mechanism
unspecified so that an agent may discover any of these ideas or a better one.
